% !TEX program = xelatex
\documentclass[12pt,a4paper]{report}

\usepackage{graphicx}
\usepackage{geometry}
\usepackage{float}
\usepackage{xcolor}
\usepackage{titlesec}
\usepackage[space]{grffile}
\usepackage{xepersian}

\geometry{top=3cm, bottom=3cm, left=2.5cm, right=3cm}

% تنظیم فاصله بعد از عنوان‌ها
\titlespacing*{\section}{0pt}{2ex plus 1ex minus .2ex}{1.5ex plus .2ex}
\titlespacing*{\subsection}{0pt}{1.8ex plus 1ex minus .2ex}{1ex plus .2ex}

% فونت فارسی ۱۴ و انگلیسی ۱۲
\settextfont[
Path = C:/USERS/SAZGAR/APPDATA/LOCAL/MICROSOFT/WINDOWS/FONTS/,
BoldFont = IRNAZANIN.TTF,
Scale = 1.166
]{IRNAZANIN.TTF}
\setlatintextfont[Scale=1]{Times New Roman}

\linespread{1.5}

\begin{document}
	\sloppy
	
	% ----------------------------------------------------------------
	% فصل ۶: طراحی و پیاده‌سازی سیستم چند عاملی حقوقی
	% ----------------------------------------------------------------
	
	\chapter*{\bfseries فصل ۶: طراحی و پیاده‌سازی سیستم چند عاملی حقوقی}
	\addcontentsline{toc}{chapter}{فصل ۶: طراحی و پیاده‌سازی سیستم چند عاملی حقوقی}
	
	\vspace{-2em}
	\noindent\rule{\textwidth}{0.5pt}
	\vspace{0.5em}
	
	% ----------------------------------------------------------------
	% بخش مقدمه فصل
	% ----------------------------------------------------------------
	\section{مقدمه}
	
	گسترش مدل‌های زبانی بزرگ در سال‌های اخیر، چشم‌انداز جدیدی برای خودکارسازی تحلیل و پاسخ‌گویی به سؤالات حقوقی گشوده است، اما اتکا به یک عامل منفرد هنوز با چالش‌های جدی در دقت، قابلیت اتکا و تبعیت سخت‌گیرانه از منابع قانونی روبه‌رو است. مطالعات حوزه «پرسش‌پاسخ حقوقی» نشان می‌دهد که دل‌های زبانی بزرگ در غیاب سازوکارهای کنترلی، مستعد تولید پاسخ‌های توهم‌آمیز و غیرقابل ردیابی هستند و صرفاً به‌کارگیری \lr{RAG} نیز همیشه برای تضمین صحت و شفافیت کفایت نمی‌کند[37, 38]. این مسئله در بافت آزمون‌های استاندارد حقوقی مانند آزمون وکالت که انتخاب نادرست یک گزینه می‌تواند نتیجه‌ای کاملاً متفاوت نسبت به منطق قانون به همراه داشته باشد، برجسته‌تر می‌شود.
	
	در پاسخ به این چالش‌ها، رویکردهای مبتنی بر چند عامل \lr{(Multi-Agent Systems)} با بهره‌گیری از چندین نقش تخصصی (مانند پژوهشگر، استدلال‌گر، منتقد و قاضی) به‌عنوان راه‌حلی نویدبخش برای افزایش دقت و پایداری استدلال حقوقی پیشنهاد شده‌اند. در حوزه حقوق، چارچوب‌هایی نظیر \lr{MALR} که از همکاری چند عامل برای تحلیل نظریه‌های حقوقی و بهبود استدلال پیچیده استفاده می‌کند، نشان داده‌اند که هماهنگی میان عامل‌های مستقل می‌تواند فهم عمیق‌ترِ مفاهیم حقوقی و استنباط‌های چندمرحله‌ای را تسهیل کند[39]. به‌طور مشابه، معماری‌های مبتنی بر مناظره و بازخورد چندعاملی برای پیش‌بینی آراء قضایی با الهام از فرایند دادرسی و مناظره در دادگاه، بهبود معناداری در دقت پیش‌بینی و کاهش خطاهای استدلالی گزارش کرده‌اند[40].
	
	افزون بر این، چارچوب‌هایی همچون \lr{AgentsCourt} و \lr{AgentsBench} با شبیه‌سازی نقش‌های مختلف دادگاه (از جمله وکلای طرفین، قاضی و عوامل بازیاب قوانین) نشان داده‌اند که ترکیب بازیابی دانش حقوقی با مناظره چندعاملی می‌تواند هم کیفیت تصمیم‌گیری و هم شفافیت و قابلیت توضیح‌پذیری سامانه‌های هوشمند حقوقی را ارتقا دهد[41, 42]. این رویکردها با جداسازی صریح مراحل «بازیابی منابع»، «استدلال مبتنی بر متن قانون» و «نقد و بازبینی» به سمت معماری‌هایی حرکت می‌کنند که نه‌تنها پاسخ نهایی، بلکه فرایند رسیدن به پاسخ را نیز قابل‌پیگیری و قابل‌داوری می‌سازند.
	
	هم‌زمان، پژوهش‌های انجام‌شده در سایر حوزه‌های تخصصی نظیر پزشکی و مالی نیز نشان داده‌اند که افزودن عامل‌های منتقد، بازتاب‌گر یا ناظر در کنار عامل استدلال‌گر اصلی می‌تواند با شناسایی و اصلاح گام‌های استدلالی نادرست، هم دقت و هم اعتمادپذیری سیستم‌های پرسش‌پاسخ را به‌طور قابل ملاحظه‌ای افزایش دهد[43, 44]. در این کارها، معماری‌های چندعاملی به‌جای اتکا به یک مدل منفرد، با توزیع مسئولیت میان چند عامل و استفاده از حلقه‌های بازخورد، توهمات مدل را کاهش داده و سازگاری پاسخ‌ها با شواهد موجود را تقویت کرده‌اند.
	
	بر این مبنا، طراحی یک سیستم چندعاملی ویژه حوزه حقوق، که در آن یک عامل پژوهشگر مسئول بازیابی هدفمند مواد قانونی، یک عامل استدلال‌گر مسئول تحلیل و انتخاب گزینه بر اساس منابع، و یک عامل منتقد/قاضی مسئول شناسایی تعارض صریح میان پاسخ و متن قانون و اعمال بازبینی باشد، پاسخی منطقی به نیازهای آزمون‌های استاندارد حقوقی محسوب می‌شود. چنین معماری‌ای با تکیه بر \lr{RAG}، زنجیره استدلال گام‌به‌گام و لایه نقد خودکار، نه‌تنها احتمال بروز پاسخ‌های توهم‌آمیز را کاهش می‌دهد، بلکه امکان کنترل کیفی پاسخ‌ها و نزدیک‌کردن رفتار مدل به منطق استنباط حقوقی انسان‌محور را نیز فراهم می‌کند.
	
	\section{معماری سیستم پیشنهادی}
	
	در این پژوهش، به‌منظور غلبه بر چالش‌های موجود در پاسخ‌دهی خودکار به سوالات حقوقی، یک معماری نوین مبتنی بر سیستم‌های چندعاملی طراحی و پیاده‌سازی شده است. این معماری با بهره‌گیری از چارچوب \lr{LangGraph}، فرآیند پاسخ‌دهی را به مجموعه‌ای از گام‌های گسسته تجزیه می‌کند که تحت نظارت یک عامل مرکزی مدیریت می‌شوند. برخلاف مدل‌های خطی سنتی (مانند زنجیره‌های ساده \lr{RAG})، این سیستم از یک گراف حالت پویا \lr{(StateGraph)} بهره می‌برد که امکان تکرار، اصلاح و بازخورد را فراهم می‌سازد[39].
	
	\subsection{ساختار کلی جریان کار}
	
	هسته اصلی سیستم، یک گراف جهت‌دار است که در آن هر گره \lr{(Node)} نمایانگر یک عامل یا واحد پردازشی مستقل و هر یال \lr{(Edge)} نشان‌دهنده جریان انتقال داده و کنترل است. در معماری پیشنهادی یک عامل ناظر مرکزی، وظیفه هماهنگی و مسیریابی بین سایر اجزا را بر عهده دارد. اجزای اصلی این گراف عبارتند از:
	
\begin{enumerate}
	\item \textbf{عامل ناظر (\lr{Supervisor Agent}):} این عامل به عنوان «مغز متفکر» سیستم عمل می‌کند. وظیفه اصلی آن تحلیل وضعیت فعلی (\lr{State}) و تصمیم‌گیری برای گام بعدی است. ناظر بر اساس اولویت‌های تعریف‌شده (مانند نیاز به بازیابی اطلاعات، نیاز به استدلال، یا نیاز به بازبینی)، کنترل را به یکی از عامل‌های تخصصی واگذار می‌کند. این رویکرد متمرکز، از سردرگمی و تداخل وظایف بین عامل‌ها جلوگیری می‌کند.
	
	\item \textbf{عامل پژوهشگر (\lr{Researcher Agent}):} مسئولیت بازیابی اسناد حقوقی مرتبط (\lr{RAG}) را بر عهده دارد. این عامل با اتصال به پایگاه دانش قوانین، مواد قانونی و آرای قضایی مرتبط با صورت سوال را استخراج و پالایش می‌کند.
	
	\item \textbf{عامل استدلال‌گر (\lr{Reasoner Agent}):} وظیفه تولید پاسخ نهایی را بر عهده دارد. این عامل با دریافت اسناد بازیابی‌شده و صورت سوال، پاسخی مستدل همراه با ارجاع دقیق به مواد قانونی تولید می‌کند.
	
	\item \textbf{عامل منتقد (\lr{Critic Agent}):} نقش «قاضی» یا «بازرس» را ایفا می‌کند. این عامل پاسخ تولیدشده توسط استدلال‌گر را با اسناد مرجع مقایسه کرده و در صورت مشاهده هرگونه تناقض صریح یا عدم انطباق با متن قانون، درخواست بازبینی (\lr{Revision}) صادر می‌کند.
	
	\item \textbf{اجراکننده ابزار (\lr{Tool Executor}):} یک گره تخصصی برای مدیریت و اجرای ابزارهای خارجی (مانند موتور جستجوی اسناد یا ابزار تحلیل گزینه‌ها) است که امکان تعامل امن مدل زبانی با محیط خارج را فراهم می‌کند.
\end{enumerate}
	
	\begin{figure}[H]
		\centering
		\includegraphics[width=0.75\textwidth]{F:/Thesis/project/3-Multi-Agent-System/data/graph.png}
		\caption{شکل ۶-۱: نمودار جریان کار سیستم چندعاملی پیشنهادی و تعاملات بین عامل‌ها تحت مدیریت عامل ناظر.}
	\end{figure}
	
	\subsection{مدیریت وضعیت}
	
	در سیستم‌های مبتنی بر \lr{LangGraph}، تمامی عامل‌ها از طریق یک حافظه مشترک به نام \lr{State} با یکدیگر تعامل دارند. در این پژوهش، از یک طرحواره اختصاصی برای مدیریت وضعیت استفاده شده است. این ساختار تضمین می‌کند که تمامی عامل‌ها به آخرین نسخه از داده‌ها دسترسی دارند و تغییرات اعمال‌شده توسط یک عامل (مثلاً اسناد بازیابی‌شده توسط پژوهشگر) بلافاصله برای سایرین (مثلاً استدلال‌گر) قابل مشاهده است.

	\section{تشریح سیستم چند عامله‌}
	
	\subsection{عامل‌ها}
	
	همانطور که گفته شد این سامانه چندعامله بر مبنای الگوی \lr{Supervisor-Orchestrated} طراحی شده است؛ بدین معنا که یک عامل ناظر مرکزی، جریان کار را در قالب یک فرایند مرحله‌ای هدایت کرده و سایر عامل‌ها را بسته به نیاز هر سؤال فعال می‌کند. برای پاسخ‌گویی به سؤالات چهارگزینه‌ای حقوقی، وظایف کلیدی به صورت تفکیک‌شده بین «بازیابی»، «استدلال»، «نقد/کنترل کیفیت» و «اجرای ابزار» توزیع می‌شود تا هم دقت افزایش یابد و هم احتمال خطاهای ناشی از توهم یا استدلال ناسازگار کاهش پیدا کند.
	
	\subsection*{عامل ناظر}
	
	عامل ناظر نقش «مدیر یا تصمیم‌گیر» را دارد و در هر گام با بررسی وضعیت مشترک سیستم، تعیین می‌کند که گام بعدی کدام است، رفتن به مرحله بازیابی، تولید پاسخ، نقد پاسخ، اجرای ابزار یا نهایی‌سازی. منطق این تصمیم‌گیری به‌صورت اولویت‌محور طراحی شده است؛ به‌طور مشخص، در صورت وجود خروجی نهایی مسیر خاتمه فعال می‌شود و در صورت وجود فراخوانی ابزار ابتدا اجرای ابزارها انجام می‌گیرد و سپس در نبود زمینه، پژوهشگر، در نبود پیش‌نویس پاسخ، استدلال‌گر، و در نبود ارزیابی منتقد، منتقد فعال می‌شود.
	
	\subsection*{عامل پژوهشگر}
	
	عامل پژوهشگر وظیفه بازیابی اسناد حقوقی مرتبط با سؤال را بر عهده دارد. خروجی پژوهشگر علاوه بر زمینه، شامل نتایج اسناد بازیابی شده و خلاصه فرادادگان نیز هست تا امکان تحلیل و ردیابی منبع فراهم شود. پژوهشگر دو شیوه اجرا دارد یکی اجرای مستقیم زنجیره بازیابی اسناد و دیگری استفاده از ابزار بازیابی از طریق ثبت یک درخواست ابزار و انتظار برای پاسخ آن. این دوحالته بودن، هم انعطاف عملیاتی ایجاد می‌کند و هم امکان کنترل دقیق‌تر روی خطاها و ثبت گزارشات سیستمی را فراهم می‌سازد.
	
	\subsection*{عامل استدلال‌گر}
	
	عامل استدلال‌گر مسئول تولید پاسخ نهایی با قالب ساختاریافته \lr{TOON} است که شامل سه مؤلفه توضیح مستدل، شماره گزینه صحیح تشخیص داده شده توسط عامل و سطح اطمینان می‌شود. در دستور متنی این عامل قید شده است که پاسخ باید «صرفاً بر اساس منابع بازیابی‌شده» تولید شود و توضیح، ارجاع دقیق به مواد/اصول منابع داشته باشد؛ همچنین خروجی باید دقیقاً در قالب \lr{TOON} ارائه گردد. در این پژوهش، فرمت TOON به‌جای JSON انتخاب شده است، زیرا TOON با ساختار ساده‌تر و حذف نشانه‌گذاری‌های زائد، خوانایی بالاتری برای انسان و درک‌پذیری بهتری برای مدل‌های زبانی فراهم می‌کند و در نتیجه موجب کاهش مصرف توکن در فرایند پردازش می‌شود. افزون بر این، انعطاف‌پذیری TOON در نمایش روابط معنایی و توضیحی، همراه با وابستگی کمتر به قواعد نحوی سخت‌گیرانه نسبت به JSON، احتمال تفسیر نادرست ساختار را کاهش داده و پایداری و یکنواختی پاسخ‌های مدل را افزایش می‌دهد.
	
	در آزمایش‌های انجام‌شده مشاهده شد که اگر پارامتر دما برای عامل استدلال‌گر از حدود \lr{۰.۲} به بالا تنظیم شود، مدل در عمل کمتر به پیشنهادها و سیگنال‌های ارائه‌شده توسط ابزار بررسی‌کننده گزینه‌ها توجه می‌کند حتی با وجود راهنمایی در ورودی، بر پاسخ تولیدی خود پافشاری نشان می‌دهد. این مشاهده با منطق طراحی سامانه هم‌راستا است، زیرا افزایش دما معمولاً باعث افزایش تنوع خروجی و کاهش «اطاعت از قیود» می‌شود؛ بنابراین برای حفظ رفتار قاعده‌مند خصوصاً در پاسخ به سوالات چهار کزینه‌ای حقوقی تنظیم دمای پایین برای عامل استدلالگر یک انتخاب طراحی مؤثر تلقی می‌گردد.
	
	\subsubsection*{نقش اعتبارسنجی گزینه‌ها در استدلال‌گر}
	
	این عامل در صورت فعال بودن ابزار اعتبارسنجی گزینه‌ها و وجود زمینه، ابتدا در صورت نبود تحلیل گزینه‌ها  یک درخواست ابزار برای تحلیل گزینه‌ها ایجاد می‌کند و سپس بعد از دریافت خروجی ابزار، آن را به‌صورت «راهنما» وارد دستور متنی می‌کند. در متن دستور متنی به‌طور صریح تأکید می‌شود که خروجی این ابزار صرفاً توصیه یا راهنمایی است و در صورت تعارض با متن منابع، باید به منابع اعتماد شود. این طراحی برای کاهش انتخاب‌های اشتباه ناشی از خطای استدلالی و نیز ایجاد یک لایه کنترل قبل از تولید پاسخ نهایی به‌کار می‌رود.
	
	یکی دیگر از یافته‌های تجربی این بود که «اعلام اشتباه حوزه سؤال» (به عنوان مثال حوزه طرح سؤال حقوق مدنی است در حالی که حقوق جزا است یا بالعکس) می‌تواند به‌طور محسوس خروجی مدل را تغییر دهد و حتی در مواردی باعث پاسخ‌گویی ناسازگار با صورت سؤال و گزینه‌ها شود. این نتیجه نشان می‌دهد که مدل‌ها به سیگنال‌های زمینه‌ای بسیار حساس‌اند و بنابراین در طراحی نهایی، از تزریق برچسب حوزه به‌عنوان ورودیِ قطعی باید با احتیاط استفاده کرد. این نکته در بخش طراحی دستور متنی و کنترل کیفیت پاسخ‌ها اهمیت مستقیم دارد.
	
	\subsection*{عامل منتقد}
	
	عامل منتقد نقش «کنترل کیفیت سخت‌گیرانه» را ایفا می‌کند و خروجی استدلال‌گر را از نظر خطاهای عینی و منبع‌محور بررسی می‌نماید. معیارهای خطا در این سامانه سه دسته‌اند: (۱) تعارض گزینه انتخابی با متن منابع، (۲) تعارض ادعاهای عامل استدلالگر با منابع، و (۳) خطای قالب‌بندی (عدم تبعیت از \lr{TOON}). منتقد صرفاً زمانی درخواست اصلاح می‌دهد که بتواند با ارجاع مشخص به منابع، تعارض صریح را نشان دهد و در غیر این صورت باید پاسخ را تأیید کند.
	
	برای کاهش خطای مثبت کاذب \lr{(false positive)} در نقد، یک «گیت حفاظتی پس از پردازش» پیاده‌سازی شده است: اگر منتقد \lr{needs\_revision=true} اعلام کند ولی در متن خروجی‌اش ارجاع معتبر به منبع وجود نداشته باشد، سیستم به‌صورت خودکار نیاز به اصلاح را \lr{false} می‌کند. این سازوکار باعث می‌شود حلقه بازبینی تنها زمانی فعال شود که منتقد واقعاً استنادپذیر و قابل‌پیگیری ایراد گرفته باشد.
	
	\subsection{ابزارها}
	
	\subsubsection*{ابزار بازیابی}
	
	ابزار بازیابی برای جستجوی اسناد حقوقی مرتبط استفاده می‌شود و توسط پژوهشگر فراخوانی می‌گردد. مزیت کلیدی این ابزار نسبت به بازیابی صرفاً متنی، پشتیبانی از متادیتا و امکان کنترل بهتر خروجی بازیابی است؛ در نتیجه، سامانه می‌تواند (در صورت نیاز) بازیابی را دقیق‌تر به سمت قانون/ماده یا نوع منبع هدایت کند و سپس \lr{context} را به شکلی استاندارد برای استدلال‌گر آماده سازد.
	
	\subsubsection*{ابزار بررسی‌کننده گزینه‌ها}
	
	این ابزار برای امتیازدهی به گزینه‌ها بر اساس منابع بازیابی‌شده طراحی شده و یک پیشنهاد به همراه میزان اطمینان و تحلیل کوتاه برای هر گزینه تولید می‌کند. نقش آن در معماری، ایجاد یک لایه «راهنمای تحلیلی» قبل از تولید پاسخ نهایی است تا انتخاب گزینه از حالت صرفاً مولد به سمت انتخابِ شواهد-محور سوق داده شود.
	
	\subsubsection*{برتری \lr{Gemini 2.5 Flash} در Option Verifier}
	
	در آزمایش‌های عملی، حساسیت عملکرد ابزار «ارزیاب گزینه‌ها» به مدل زبانی زیربنایی آن به‌وضوح مشاهده شد. نتایج نشان داد که کیفیت این ابزار صرفاً به صحت بازیابی اسناد وابسته نیست، بلکه به توانایی مدل در «تحلیل دقیق متن قانونی» و نگاشت منطق آن به گزینه‌ها نیز بستگی حیاتی دارد. این یافته بر اهمیت انتخاب مدل مناسب برای ابزارهای تحلیلی در سامانه‌های حقوقی تأکید می‌کند، زیرا خطای یک ابزار می‌تواند کل زنجیره استدلال را منحرف سازد.
	
	برای نمونه، در ارزیابی سؤالی پیچیده پیرامون استثنائات تبدیل مجازات‌های تعزیری، با وجود بازیابی صحیح ماده ۳۷ قانون مجازات اسلامی، اغلب مدل‌ها در تحلیل متن و انتخاب گزینه صحیح دچار خطا شدند، در حالی که مدل \lr{Gemini 2.5 Flash} توانست با تحلیل دقیق‌تری، پیشنهاد درستی ارائه دهد. جزئیات این نمونه در جدول زیر ارائه شده است.
	
	\begin{table}[H]
		\centering
		\begin{tabular}{p{0.18\textwidth} p{0.72\textwidth}}
			\hline
			\textbf{سؤال} &
			کدام مجازات تعزیری زیر در مقام تخفیف قابل تبدیل به مجازات دیگر نیست؟\\
			& ۱) مصادره اموال \\
			& ۲) محرومیت از حقوق اجتماعی \\
			& ۳) حبس درجه ۳ \\
			& ۴) شلاق درجه ۶ \\ \hline
			\textbf{ماده ۳۷ قانون مجازات اسلامی} &
			\lr{}در صورت وجود یک یا چند جهت از جهات تخفیف، دادگاه می‌تواند مجازات تعزیری را به‌ نحوی که به حال متهم مناسب‌تر باشد به‌شرح زیر تقلیل دهد یا تبدیل کند: \\
			& الف ـ تقلیل مجازات حبس به میزان یک تا سه درجه در مجازات‌های درجه چهار و بالاتر. \\
			& ب ـ تقلیل مجازات حبس درجه پنج و درجه شش به میزان یک تا دو درجه یا تبدیل این مجازات و مجازات حبس درجه هفت حسب مورد به جزای نقدی متناسب با همان درجه. \\
			& پ ـ تبدیل مصادره کل اموال به جزای نقدی درجه یک تا چهار. \\
			& ت ـ تقلیل انفصال دائم به انفصال موقت به میزان پنج تا پانزده سال. \\
			& ث ـ تقلیل سایر مجازات‌های تعزیری به میزان یک یا دو درجه یا تبدیل آن به مجازات دیگر از همان درجه یا یک درجه پایین‌تر. \\
			& تبصره: چنانچه در اجرای مقررات این ماده یا سایر مقرراتی که به‌موجب آن مجازات تخفیف می‌یابد، حکم به حبس کمتر از نود و یک روز صادر شود، به مجازات جایگزین مربوط تبدیل می‌شود. \\ \hline
		\end{tabular}
		\caption{سوال 86 آزمون وکالت 1403 همراه با ماده قانونی مرتبط}
		\label{tab:option-verifier-article37}
	\end{table}
	
	
	در طراحی اولیه، دستور متنی همه اجزا به زبان انگلیسی تنظیم شده بود، اما در ادامه مشاهده شد که برای \lr{Option Verifier}، تغییر زبان دستور متنی به فارسی باعث بهبود ملموس در درک متون حقوقی فارسی و تشخیص گزینه صحیح شد. این تغییر منجر به بهبود ملموس در توانایی این ابزار برای درک ظرایف متون حقوقی فارسی و تشخیص گزینه صحیح شد. تحلیل این نتیجه نشان می‌دهد که وظیفه \lr{Option Verifier}، یعنی مقایسه تحلیلی و دقیق گزینه‌ها با متن پیچیده و فنی قوانین، بیش از سایر اجزا به درک عمیق معنایی و اصطلاحی زبان منبع وابسته است. استفاده از دستور متنی فارسی در این ابزار خاص، بار شناختی مدل برای ترجمه و تطبیق مفاهیم را از بین برده و به آن اجازه می‌دهد تمام توان پردازشی خود را بر تحلیل حقوقی متمرکز کند. این یافته یک توصیه طراحی مهم را به همراه دارد، در معماری‌های چندزبانه، برای ابزارها یا عامل‌هایی که وظیفه تحلیل معنایی عمیق بر روی متون زبان مقصد را بر عهده دارند، استفاده از دستور متنی هم‌زبان با منبع می‌تواند یک عامل کلیدی در افزایش دقت و کاهش خطاهای تفسیری باشد.
	
	\subsection{مکانیزم بازخورد و اصلاح}
	
	یکی از نوآوری‌های کلیدی این معماری، پیاده‌سازی یک حلقه بازخورد و اصلاح تکرارشونده با الهام از فرآیندهای نقد و همتاسنجی \lr{(Peer Review)} است[40, 45]. این مکانیزم، که از مفاهیمی چون «خوداصلاحی» (\lr{Self-Correction}) و «مناظره محاسباتی» (\lr{Computational Debate}) در سیستم‌های هوشمند الهام گرفته است، با هدف افزایش دقت و پایبندی پاسخ به منابع حقوقی طراحی شده است. برخلاف زنجیره‌های خطی که در آن پاسخ در یک گام تولید و نهایی می‌شود، در این سیستم پاسخ اولیه به‌عنوان یک «پیش‌نویس» در نظر گرفته شده و تا زمان کسب تأیید از عامل منتقد، در یک چرخه اصلاحی قرار می‌گیرد.
	
	این چرخه به‌صورت یک تعامل پویا بین دو عامل استدلال‌گر و منتقد، تحت نظارت عامل ناظر، به شرح زیر عمل می‌کند:
	
	\begin{enumerate}
		\item \textbf{تولید پیش‌نویس:} در گام نخست، عامل استدلال‌گر بر اساس اسناد بازیابی‌شده، یک پاسخ اولیه در قالب \lr{draft\_toon} تولید می‌کند.
		
		\item \textbf{ارسال برای نقد:} عامل ناظر، با تشخیص وجود \lr{draft\_toon} و عدم وجود \lr{critic\_toon}، کنترل را به عامل منتقد واگذار می‌کند.
		
		\item \textbf{ارزیابی منتقد:} عامل منتقد، پیش‌نویس را با اسناد مرجع مقایسه کرده و در صورت مشاهده هرگونه تناقض صریح، یک خروجی \lr{critic\_toon} با وضعیت  "نیاز به بررسی" تولید می‌کند. این خروجی صرفاً یک رد یا تأیید ساده نیست، بلکه شامل توضیح علت خطا و اقدام اصلاحی پیشنهادی نیز هست تا به استدلال‌گر در فرآیند اصلاح کمک کند.
		
		\item \textbf{حلقه اصلاح:}
		\begin{itemize}
			\item \textbf{الف)} در صورت نیاز به اصلاح اگر شمارنده بازبینی از سقف مجاز فراتر نرفته باشد، ناظر مجدداً کنترل را به عامل استدلال‌گر بازمی‌گرداند. در این مرحله، دستور متنی استدلال‌گر به‌صورت پویا با افزودن بازخورد منتقد به‌روزرسانی می‌شود و از آن خواسته می‌شود تا خطای مشخص‌شده را برطرف سازد. پس از تولید پیش‌نویس جدید، شمارنده بازبینی یک واحد افزایش یافته و چرخه از گام ۲ تکرار می‌شود.
			
			\item \textbf{ب)} در صورت تأیید اگر پاسخ تأیید شود و نیاز به بازبینی نداشته باشیم، ناظر چرخه را شکسته و کنترل را به گره نهایی‌سازی منتقل می‌کند.
		\end{itemize}
	\end{enumerate}
	
	این فرآیند تکرارشونده تضمین می‌کند که پاسخ نهایی، محصول یک استدلال یک‌باره و کنترل‌نشده نیست، بلکه حاصل یک فرآیند پالایش است که در آن خطاهای احتمالی شناسایی و اصلاح شده‌اند. تعیین یک سقف برای تعداد بازبینی‌ها نیز یک مکانیزم حفاظتی ضروری است که از گیر افتادن سیستم در حلقه‌های اصلاح بی‌حاصل جلوگیری کرده و خاتمه یافتن فرآیند را در زمان معقول تضمین می‌کند. این معماری نقد-محور، سیستم را قادر می‌سازد تا به سطح بالاتری از اعتمادپذیری و سازگاری با منابع دست یابد که در حوزه حساس حقوق امری ضروری است[42].

	\section{ارزیابی و نتایج تجربی}
	
	در این بخش، عملکرد سامانه چندعاملی پیشنهادی مورد ارزیابی دقیق قرار گرفته و نتایج آن با رویکرد پایه (\lr{RAG} استاندارد) مقایسه شده است. هدف از این آزمایش‌ها، سنجش تجربی این فرضیه است که آیا تفکیک وظایف در قالب یک معماری \lr{Supervisor-Orchestrated} و افزودن لایه‌های نقد و تحلیل گزینه، می‌تواند بر محدودیت‌های ذاتی مدل‌های زبانی در استدلال حقوقی غلبه کند یا خیر. نتایج به‌دست‌آمده نه‌تنها میزان دقت کلی، بلکه رفتار مدل‌ها در زیرشاخه‌های مختلف حقوقی را نیز آشکار می‌سازد.
	
	\subsection{تنظیمات آزمایش}
	
	برای انجام این ارزیابی، چهار مدل زبانی بزرگ که در مراحل غربالگری اولیه از میان ۱۴ مدل نامزد، بیشترین دقت، سرعت و سازگاری را با زبان فارسی و متون حقوقی نشان داده بودند، انتخاب شدند. این مدل‌ها شامل \lr{Gemini 2.5 Flash} (بهترین مدل تجاری و بهترین مدل از نظر دقت)، \lr{Llama 4 Maverick} ،\lr{Qwen3 80} (بهترین مدل متن باز) و \lr{Qwen3 235} هستند. مجموعه داده مورد استفاده برای آزمون، شامل ۱۲۰ سوال چهارگزینه‌ای از آزمون وکالت سال ۱۴۰۳ است که به دلیل ماهیت استدلالی و پیچیدگی بالا، معیار مناسبی برای سنجش توانایی مدل‌ها در حل مسائل واقعی حقوقی محسوب می‌شود. تمامی آزمایش‌ها در یک محیط یکسان و با پارامترهای ثابت اجرا شدند؛ به طور مشخص، برای عامل استدلال‌گر دمای تولید روی \lr{۰.۱} تنظیم شد تا پاسخ‌ها قطعی و تکرارپذیر باشند و سقف تعداد بازبینی مجاز نیز برابر با ۲ در نظر گرفته شد.
	
	\subsection{نتایج کلی}
	
	بررسی نتایج کلی نشان‌دهنده یک جهش کیفی در عملکرد مدل‌ها پس از استقرار در معماری چندعاملی است. مقایسه داده‌های حاصل از این آزمایش با نتایج فصل پیشین (\lr{RAG} ساده) حاکی از آن است که تمامی چهار مدل منتخب، در بستر سیستم چندعاملی دقت بالاتری را ثبت کرده‌اند. برجسته‌ترین بهبود عملکرد مربوط به مدل \lr{Llama 4 Maverick} است که دقت آن از \lr{۵۷.۹۸٪} در حالت \lr{RAG} به \lr{۶۴.۵٪} در حالت چندعاملی افزایش یافته است. این رشد \lr{۶.۵٪}، به‌وضوح نشان می‌دهد که این مدل، که ذاتاً از توانایی استدلالی خوبی برخوردار است، در حالت \lr{RAG} ساده دچار خطاهای ناشی از عدم تمرکز یا تفسیر شتاب‌زده بوده است؛ اما با قرارگیری در چرخه نقد و اصلاح (\lr{Critic-Refinement Loop})، توانسته است استدلال‌های خود را تصحیح کرده و به پاسخ صحیح نزدیک‌تر شود.
	
	مدل \lr{Qwen3 235} نیز الگوی مشابهی را با افزایش دقت از \lr{۵۷.۱۴٪} به \lr{۶۲.۸٪} تجربه کرده است که نشان‌دهنده تأثیر مثبت ابزار تحلیل گزینه در هدایت مدل‌های مقیاس‌بزرگ است. در این میان، مدل \lr{Gemini 2.5 Flash} با ثبت دقت \lr{۶۴.۷٪}، بالاترین عملکرد کلی را در میان تمامی مدل‌ها به خود اختصاص داد. اگرچه میزان بهبود این مدل نسبت به حالت پایه (\lr{۶۳.۸۷٪}) کمتر از سایرین است، اما پایداری بالای آن نشان می‌دهد که این مدل حتی در حالت پایه نیز خطاهای فاحش کمتری داشته و معماری چندعاملی بیشتر نقش تثبیت‌کننده را برای آن ایفا کرده است. در مجموع، این یافته‌ها اثبات می‌کند که مکانیزم‌های کنترلی تعبیه‌شده در سیستم، نظیر بررسی مجدد توسط منتقد و تحلیل جداگانه گزینه‌ها، تأثیر معناداری در کاهش توهمات و افزایش قابلیت اتکای سیستم داشته‌اند[40].
	
	\begin{figure}[H]
		\centering
		\includegraphics[width=0.95\textwidth]{F:/Thesis/project/3-Multi-Agent-System/data/model acc comparison.png}
		\caption{شکل ۶-۲: نمودار مقایسه عملکرد کلی مدل‌ها در دو حالت \lr{RAG} ساده و سیستم چندعاملی}
	\end{figure}
	
	\subsection{تحلیل تفکیکی بر اساس دروس}
	
	تحلیل عملکرد مدل‌ها به تفکیک دروس حقوقی، تصویری دقیق‌تر از نقاط قوت و ضعف هر معماری ارائه می‌دهد و نشان می‌دهد که «هوش حقوقی» مدل‌ها در تمام حوزه‌ها یکسان توزیع نشده است. در درس حقوق اساسی، تمامی مدل‌ها عملکردی درخشان و بالاتر از میانگین داشته‌اند؛ به طوری که \lr{Llama 4 Maverick} با دستیابی به دقت خیره‌کننده \lr{۸۴.۲۱٪}، رکوردی قابل توجه ثبت کرده است. این موفقیت را می‌توان به ماهیت اصول قانون اساسی نسبت داد که معمولاً صریح‌تر هستند و کمتر از مباحثی مانند حقوق مدنی نیاز به ترکیب چند ماده قانونی دارند.
	
	در نقطه مقابل، درس حقوق جزا به عنوان چالش‌برانگیزترین بخش برای مدل‌های خانواده \lr{Qwen} شناسایی شد. مدل \lr{Qwen3 235} در این درس با افتی شدید، تنها موفق به کسب دقت \lr{۳۱.۵۷٪} گردید، در حالی که مدل‌های \lr{Gemini} و \lr{Llama} توانستند دقت خود را در محدوده قابل قبول (بالای \lr{۵۰٪}) حفظ کنند. این تفاوت فاحش نشان می‌دهد که مدل‌های \lr{Qwen} احتمالاً در درک ظرایف زبانی مرتبط با عناصر جرم و مجازات‌ها در زبان فارسی با دشواری مواجه هستند و صرفاً افزایش پارامترهای مدل نتوانسته است این ضعف زبانی-مفهومی را جبران کند.
	
	یکی از یافته‌های جالب و غیرمنتظره این پژوهش، عملکرد مدل \lr{Qwen3 80} در درس حقوق مدنی بود. این مدل با ثبت دقت \lr{۶۹.۲۳٪}، نه‌تنها از برادر بزرگتر خود (\lr{Qwen3 235}) پیشی گرفت، بلکه عملکردی بهتر از مدل قدرتمند \lr{Gemini} (با \lr{۵۵٪}) ارائه داد. از آنجا که سوالات حقوق مدنی (به‌ویژه مباحث ارث و قراردادها) نیازمند استدلال‌های منطقی و شبه‌ریاضی دقیق هستند، این نتیجه نشان می‌دهد که مدل \lr{Qwen3 80} احتمالاً دارای توانایی استنتاج منطقی (\lr{Logical Reasoning}) بالاتری است، هرچند که ممکن است در درک متون توصیفی ضعیف‌تر عمل کند.
	
	در نهایت، بررسی درس‌های آیین دادرسی کیفری و حقوق تجارت نشان‌دهنده پایداری و تعادل عملکرد در مدل \lr{Llama 4 Maverick} است که به ترتیب دقت \lr{۷۵٪} و \lr{۷۰٪} را ثبت کرده است. این توازن عملکرد در دروس مختلف، \lr{Llama 4} را به عنوان متعادل‌ترین مدل برای کاربردهای عمومی حقوقی معرفی می‌کند، در حالی که مدل‌هایی مانند \lr{Qwen3 80} می‌توانند به عنوان متخصصین حوزه‌های خاص (مانند مدنی) در نظر گرفته شوند. این ناهمگونی در نتایج دروس مختلف، ایده استفاده از رویکرد ترکیبی (\lr{Ensemble}) را تقویت می‌کند؛ سیستمی که بتواند برای سوالات مدنی به \lr{Qwen} و برای سوالات اساسی به \lr{Llama} ارجاع دهد، پتانسیل دستیابی به دقتی فراتر از هر مدل منفرد را خواهد داشت.
	
	\begin{figure}[H]
		\centering
		\includegraphics[width=0.95\textwidth]{F:/Thesis/project/3-Multi-Agent-System/data/acc by cat.png}
		\caption{شکل ۶-۳: نمودار عملکرد تفکیکی مدل‌ها در دروس شش‌گانه حقوقی (سیستم چندعاملی)}
	\end{figure}
	
	\section{بحث و تحلیل}
	
	نتایج تجربی به‌دست‌آمده، اثربخشی معماری چندعاملی را در بهبود دقت پاسخ‌دهی به سؤالات حقوقی تأیید می‌کند. با این حال، درک عمیق دلایل این بهبود و همچنین شناسایی چالش‌ها و محدودیت‌های ذاتی این رویکرد، برای ارزیابی جامع پژوهش ضروری است. این بخش به تحلیل کیفی این دو جنبه می‌پردازد.
	
	\subsection{اثر تفکیک وظایف بر کاهش خطا}
	
	افزایش دقت مشاهده‌شده در مدل‌ها، بیش از آنکه حاصل توانمندی یک مدل خاص باشد، نتیجه‌ی مستقیم «تفکیک وظایف» در معماری سیستم است. در رویکردهای تک‌مرحله‌ای مانند \lr{RAG} ساده، یک مدل زبانی واحد موظف است به‌صورت همزمان چندین وظیفه شناختی پیچیده را انجام دهد، خواندن و درک سؤال، بازیابی اسناد مرتبط، تحلیل منابع، ارزیابی هر چهار گزینه، استنتاج پاسخ نهایی و تدوین توضیح مستدل. این حجم از وظایف احتمال بروز خطا، تفسیر شتاب‌زده و نادیده گرفتن جزئیات کلیدی را افزایش می‌دهد.
		
	بنابراین، معماری چندعاملی با تبدیل یک مسئله پیچیده به زنجیره‌ای از مسائل ساده‌تر و قابل‌کنترل، یک چارچوب قاعده‌مند برای استدلال ایجاد می‌کند که در آن هر عامل، خروجی عامل قبلی را کنترل و پالایش می‌کند و احتمال رسیدن خطای اولیه به خروجی نهایی را کاهش می‌دهد.
	
	\subsection{چالش‌ها و محدودیت‌ها}
	
	علی‌رغم مزایای چشمگیر، معماری چندعاملی با چالش‌ها و محدودیت‌های عملی نیز همراه است که باید مورد توجه قرار گیرند:
	
	\textbf{افزایش تأخیر و هزینه محاسباتی:} بارزترین محدودیت این سیستم، افزایش زمان پاسخ‌دهی است. هر گام در گراف (فراخوانی پژوهشگر، استدلال‌گر و منتقد) معادل یک یا چند فراخوانی به مدل زبانی است. یک چرخه کامل اصلاح (استدلال‌گر \lr{←} منتقد \lr{←} استدلال‌گر) حداقل سه فراخوانی مدل‌ربانی بزرگ را در بر می‌گیرد، در حالی که در \lr{RAG} ساده تنها یک فراخوانی وجود دارد. این سربار محاسباتی (\lr{Computational Overhead}) نه‌تنها منجر به کندی سیستم می‌شود، بلکه هزینه مالی ناشی از مصرف توکن را نیز به شکل قابل توجهی افزایش می‌دهد[40, 41].
	
	\textbf{پیچیدگی پیاده‌سازی و شکنندگی سیستم:} مدیریت یک سیستم چندعاملی، هماهنگ‌سازی حالت‌ها و اطمینان از صحت تعاملات بین عامل‌ها، بسیار پیچیده‌تر از اجرای یک زنجیره خطی است. سیستم به دلیل وابستگی به فرمت‌های دقیق خروجی می‌تواند شکننده باشد. یک خطای جزئی در قالب‌بندی خروجی یک عامل می‌تواند باعث شکست کل زنجیره شود. هرچند مکانیزم‌های حفاظتی (مانند \lr{Post-check Gate} در منتقد) برای کاهش این شکنندگی تعبیه شده‌اند، اما این خود نشان‌دهنده پیچیدگی ذاتی نگهداری چنین سیستمی است.
	
	\textbf{خطر انتشار خطا:} این معماری بر پایه این فرض استوار است که هر عامل وظیفه خود را به درستی انجام می‌دهد. اما اگر یک عامل در مراحل اولیه دچار خطا شود، این خطا می‌تواند به مراحل بعدی انتشار یابد. برای مثال، اگر عامل پژوهشگر در بازیابی ماده قانونی صحیح شکست بخورد، حتی بهترین استدلال‌گر و منتقد نیز قادر به تولید پاسخ صحیح نخواهند بود، زیرا بر مبنای اطلاعات غلط استدلال می‌کنند. به همین ترتیب، یک نقد اشتباه از سوی منتقد می‌تواند استدلال‌گر را به سمت اصلاح یک پاسخ صحیح و جایگزینی آن با پاسخی نادرست سوق دهد.
	
	\section{نتیجه‌گیری فصل}
	
	در این فصل، یک معماری نوین چندعاملی برای پاسخ‌دهی به سؤالات استاندارد حقوقی طراحی، پیاده‌سازی و ارزیابی شد. هدف اصلی این پژوهش، غلبه بر محدودیت‌های سامانه‌های تک‌عاملی و افزایش دقت و اعتمادپذیری پاسخ‌ها از طریق تفکیک وظایف و ایجاد یک حلقه بازخورد انتقادی بود. سیستم پیشنهادی، مبتنی بر الگوی \lr{Supervisor-Orchestrated}، با بهره‌گیری از چهار ایجنت تخصصی (ناظر، پژوهشگر، استدلال‌گر و منتقد)، فرآیند پیچیده استدلال حقوقی را به گام‌های کوچک‌تر و قابل‌کنترل تجزیه کرد.
	
	نتایج تجربی به‌وضوح نشان داد که این معماری تأثیر مثبت و معناداری بر عملکرد تمامی مدل‌های زبانی منتخب داشته است. افزایش دقت در تمام مدل‌ها، به‌ویژه رشد چشمگیر \lr{۶.۵٪} در مدل \lr{Llama 4 Maverick}، اثربخشی مکانیزم بازخورد و اصلاح \lr{(Iterative Refinement Loop)} را در شناسایی و رفع خطاهای پنهان به اثبات رساند. این سیستم با دستیابی به دقت \lr{۶۴.۷٪}، یک معیار عملکرد جدید برای این مجموعه داده ثبت کرد. تحلیل تفکیکی نتایج بر اساس دروس حقوقی نیز نشان داد که هیچ مدلی در تمامی حوزه‌ها برتر نیست و توانایی‌های استدلالی مدل‌ها بسته به ماهیت دروس (تفسیری، منطقی یا محاسباتی) متفاوت است.
	
	با وجود این موفقیت‌ها، این رویکرد با چالش‌هایی نظیر افزایش تأخیر زمانی و هزینه محاسباتی نیز همراه است که یک بده‌بستان ذاتی بین دقت و کارایی را به نمایش می‌گذارد. با این حال، این پژوهش نشان داد که حرکت از مدل‌های یکپارچه \lr{(Monolithic)} به سمت معماری‌های چندعاملی و مبتنی بر همکاری، گامی مؤثر در جهت ساخت سیستم‌های هوشمند حقوقی «قابل اعتمادتر» و «سازگارتر با شواهد» است. چارچوب ارائه‌شده نه‌تنها دقت پاسخ نهایی را بهبود بخشید، بلکه با شفاف‌سازی فرآیند استدلال و نقد، مسیری به‌سوی هوش مصنوعی حقوقی توضیح‌پذیر \lr{(Explainable AI)} را نیز هموار می‌سازد.
	
\end{document}
